\documentclass[10pt,a4paper]{amsart}
\usepackage[margin=19mm]{geometry}
\usepackage{amsmath,amssymb,mathtools}
\usepackage[scr=rsfs,cal=euler]{mathalfa}
\usepackage{xfrac}
\usepackage[unicode,hidelinks,bookmarks=false]{hyperref}
\newcommand{\ie}{i.e.\ }
\newcommand{\cf}{cf.\ }
\newcommand{\resp}{respectively\ }
\newcommand{\Subsection}{\subsection}
\providecommand{\nobreakdash}{\nobreak\hbox{-}}
\setlength{\emergencystretch}{2em}
\multlinegap0pt
\usepackage{bm}
\usepackage{bbm}
\usepackage{dsfont}
\usepackage{xspace}
\usepackage{cancel}
\usepackage{mathtools}
\usepackage{amsthm}
\makeatletter
\@ifundefined{theorem}{\newtheorem{theorem}{Theorem}}{}
\@ifundefined{assumption}{\newtheorem{assumption}[theorem]{Assumption}}{}
\@ifundefined{lemma}{\newtheorem{lemma}[theorem]{Lemma}}{}
\makeatother
\DeclarePairedDelimiter{\abs}{|}{|}
\DeclarePairedDelimiter{\braces}{\lbrace}{\rbrace}
\DeclarePairedDelimiter{\paren}{(}{)}
\newcommand{\E}{\mathbb{E}}
\newcommand{\MSP}{\textsf{MSP}}
\newcommand{\R}{\mathbb{R}}
\newcommand{\alg}{\mathcal{A}}
\newcommand{\bX}{\mathbf{X}}
\newcommand{\binspace}{\braces{\pm 1}^\nfeats}
\newcommand{\bx}{\mathbf{x}}
\newcommand{\cellfeatindices}{J}
\newcommand{\cell}{\mathcal{C}}
\newcommand{\coordindices}{[\nfeats]}
\newcommand{\data}[1][n]{\mathcal{D}_{#1}}
\newcommand{\fcoefs}{\mathcal{S}}
\newcommand{\monomial}{\chi}
\newcommand{\nfeats}{d}
\newcommand{\nsamples}{n}
\newcommand{\querypath}{\mathcal{Q}}
\newcommand{\sparsity}{s}
\newcommand{\suppset}{S^*}
\newcommand{\traversalsize}{\sparsity_{\operatorname{T}}}
\newcommand{\traversal}{T}
\newcommand{\xmeasure}{\nu}
\renewcommand{\P}{\mathbb{P}}
\renewcommand{\xspace}{\mathcal{X}}
\title[Statistical-Computational Trade-offs for Recursive Adaptive ]{Statistical-Computational Trade-offs for Recursive Adaptive Partitioning Estimators}
\author{Yan Shuo Tan, Jason M. Klusowski, Krishnakumar Balasubramanian}
\date{Source: arXiv:2411.04394v3 (CC BY 4.0)}
\begin{document}
\maketitle

\noindent\emph{Source and licence.} Statistical-Computational Trade-offs for Recursive Adaptive Partitioning Estimators. Version arXiv:2411.04394v3 (CC BY 4.0), distributed under the Creative Commons Attribution 4.0 International licence, as recorded in the arXiv licence metadata for this exact version and shown on the abstract page https://arxiv.org/abs/2411.04394v3. Full rights notice: licenses/arxiv-2411.04394-CC-BY-4.0.txt.\par\smallskip

\noindent\textit{Editorial scope.} Counted unit: the Lemma whose source label is \texttt{lem:nonmsp\_feature\_selection\_prob} in the article (source0.tex lines 2323--2330, source label \texttt{lem:nonmsp\_feature\_selection\_prob}), delivered together with the article's own complete original proof of it (source0.tex lines 2332--2403, one \texttt{proof} environment of 72 lines). No mathematics is written, repaired or re-derived: statement and proof are the article's own text line for line. Required context retained uncounted: eq:nonparametric\_regression\_model (lines 370--373); assumption (lines 1051--1060); lem:bound\_for\_tree\_depth (lines 2278--2291). Every definition, display and label of those lines is retained unchanged. Omitted from this package and not redistributed: 1--369, 374--1050, 1061--2277, 2292--2322, 2331--2331, 2404--3788. Nothing inside a retained excerpt is omitted or altered: every retained body is delivered exactly as the article's lines are written, so no omission receipt of any kind accompanies this package. Citations: the retained excerpts carry no citation command at all, so no bibliography entry of the article is transcribed and no reference is redistributed. Reference adaptation: none was needed and none was made; every reference inside the delivered unit resolves inside it, so no unresolved reference or citation is produced. Printed slips kept literal: no printed word, symbol, hypothesis or number of the counted statement or of its proof was altered. Layout and class: the article's own class is \texttt{article} with packages including natbib, authblk, amsmath, amssymb, amsthm, mathtools, inputenc, geometry, times, hyperref, url, multirow, xr, fontawesome; the wrapper is the lane's pinned \texttt{amsart} editorial wrapper at 10pt on A4, which restates the theorem-like environments the retained text needs and adds the article's own macro definitions that the retained text needs (\texttt{\textbackslash E}, \texttt{\textbackslash MSP}, \texttt{\textbackslash P}, \texttt{\textbackslash R}, \texttt{\textbackslash abs}, \texttt{\textbackslash alg}, \texttt{\textbackslash bX}, \texttt{\textbackslash binspace}, \texttt{\textbackslash braces}, \texttt{\textbackslash bx}, \texttt{\textbackslash cell}, \texttt{\textbackslash cellfeatindices}, \texttt{\textbackslash coordindices}, \texttt{\textbackslash data}, \texttt{\textbackslash fcoefs}, \texttt{\textbackslash monomial}, \texttt{\textbackslash nfeats}, \texttt{\textbackslash nsamples}, \texttt{\textbackslash paren}, \texttt{\textbackslash querypath}, \texttt{\textbackslash sparsity}, \texttt{\textbackslash suppset}, \texttt{\textbackslash traversal}, \texttt{\textbackslash traversalsize}, \texttt{\textbackslash xmeasure}, \texttt{\textbackslash xspace}), with extra wrapper packages (bm, bbm, dsfont, xspace, cancel, mathtools). The delivered numbering is the unit's own. Assets: no retained span contains a figure environment, a graphics-inclusion command, a PDF-page inclusion, a table or a TikZ picture; no image file is copied into this package, the package declares no asset dependency and no image file is redistributed with it. Licence: version arXiv:2411.04394v3 of this article is distributed under the Creative Commons Attribution 4.0 International licence, as recorded in the arXiv licence metadata for this exact version and shown on the abstract page https://arxiv.org/abs/2411.04394v3. A full rights notice is delivered at licenses/arxiv-2411.04394-CC-BY-4.0.txt. Characters: every retained excerpt is pure ASCII, so the delivered engine, XeLaTeX with its default Latin Modern text font, reports no missing character.
\begin{equation}
    \label{eq:nonparametric_regression_model}
    Y_i = f^*(\bX_i) + \varepsilon_i, \quad i =1,2,\ldots,\nsamples.
\end{equation}

\begin{assumption}
    \label{assumption}
    Assume a binary covariate space $\xspace = \binspace$ with $\xmeasure$ the uniform measure.
    Let $f^*_0\colon \lbrace\pm 1\rbrace^\sparsity \to \R$ be any function.
    Let $\suppset \subset \coordindices$ be a subset of size $\abs*{\suppset} = \sparsity$.
    % that is drawn uniformly at random.
    We consider the nonparametric model \eqref{eq:nonparametric_regression_model} with the regression function $f^*\colon \binspace \to \R$ defined via $f^*(\bx) = f^*_0(\bx^{\suppset})$.
    % We denote the class of all such functions via $\sparsefuncs$.
    % and that $f^*$ depends only on $\sparsity$ covariates, i.e., $f^*(\bx) = h(\bx_{\suppset})$ for some $\suppset \subset \coordindices$ with $\abs*{S} \leq \sparsity$, and $h\colon \R^{|\suppset|} \to \R$.
\end{assumption}

\begin{lemma}
% [Upper bound for tree depth]
    \label{lem:bound_for_tree_depth}
    For any nonnegative integer $t$, we have the tail bound
    \begin{equation}
    \label{eq:tree_depth_probability}
        \P\braces*{\abs*{J(\bx;\data,\Theta)} \geq t} \leq \min\braces*{\frac{n}{2^{t}},1}.
    \end{equation}
    Furthermore, we have the expectation bound
    \begin{equation}
    \label{eq:tree_depth_expectation}
        \E\braces*{\abs*{J(\bx;\data,\Theta)}} \leq \log_2 \nsamples + 2.
    \end{equation}
\end{lemma}

\begin{lemma}
    \label{lem:nonmsp_feature_selection_prob}
    Under Assumption \ref{assumption}, let $f_0^*$ be a function not satisfying \MSP, and suppose $\alg_t$ is greedy.
    Let $\traversal \subset \coordindices\backslash\suppset_{\MSP}$ be a traversal for $\fcoefs\backslash\fcoefs_{\MSP}$ of size $\traversalsize$ and let $\sparsity_{\MSP} \coloneqq \abs*{\suppset_{\MSP}}$.
    % For any $0 < \delta < 1$,
    % suppose $\nsamples \leq 2^{\nfeats/\sparsity\delta}$.
    Then for each fixed query point $\bx \in \binspace$, support set $\suppset$, and random seed $\Theta$, the probability of selecting any covariate in the traversal along the query path $\querypath(\bx;\data,\Theta)$ is upper bounded as $\P_{\data}\braces*{J(\bx;\data,\Theta)\cap \traversal \neq \emptyset } \leq \frac{\traversalsize(\log_2 \nsamples + 2)}{\nfeats - \sparsity_{\MSP}-\traversalsize +1}$.
\end{lemma}

\begin{proof}
    Without loss of generality, assume $\suppset_{\MSP} = \braces*{1,2,\ldots,\sparsity_{\MSP}}$ and
    \begin{equation}
        \traversal = \braces*{\sparsity_{\MSP}+1,\sparsity_{\MSP}+1,\ldots,\sparsity_{\MSP}+\traversalsize}.
    \end{equation}
    \textit{Step 1: Defining the path coupling.}
    For $a=1,2,\ldots,\traversalsize$, define the subset $\traversal^{(a)} \coloneqq \traversal\backslash\braces*{\sparsity_{\MSP}+a}$ and the modified dataset $\data^{(a)} \coloneqq \braces*{(\bX_{i,-\traversal^{(a)}},Y_i)}_{i=1}^\nsamples$, i.e., so that all covariates with indices in $\traversal^{(a)}$ are dropped from each observation.
    We will analyze the query paths, $\querypath(\bx;\data^{(a)},\Theta)$ obtained from these modified datasets and compare them with the original.
    % For the sake of notational convenience, we also denote $\data^{(0)} \coloneqq \data$.
    % Next, for $a=0,1,2$, denote $\path^{(a)} \coloneqq \path\paren*{\bx_0,\data^{(a)},\Theta}$.
    % We further use superscripts to denote all quantities related to these query paths.
    % In particular, we let $l^{(a)}$ denote the length of $\path^{(a)} \eqqcolon \paren*{\cell_0^{(a)},\cell_1^{(a)},\ldots,\cell_{l^{(a)}}^{(a)}}$ and $\featindices_\infty^{(a)} \coloneqq \paren*{j_1^{(a)},j_2^{(a)},\ldots,j_{l^{(a)}}^{(a)}}$ denote the sequence of features split on in $\path^{(a)}$.

    \textit{Step 2: Coupled paths select covariates randomly.}
    For $a=1,2,\ldots,\traversalsize$, we claim that each response $Y_i$ is independent of the covariate vector $\bX_{i,-\traversal^{(a)}\cup\suppset_{\MSP}}$.
    To see this, consider $S \in \fcoefs$.
    If $S \in \fcoefs_{\MSP}$, then $\monomial_S$ does not depend on $\bX_{i,-\traversal^{(a)}\cup\suppset_{\MSP}}$ and is hence trivially independent.
    Otherwise, $S \in \fcoefs\backslash\fcoefs_{\MSP}$, and by the definition of a traversal, we have
    \begin{equation}
        \abs*{S\cap \traversal^{(a)}} \geq \abs*{S\cap\traversal} - 1 \geq 1.
    \end{equation}
    As such, for any fixed value of $\bx_{i,-\traversal^{(a)}\cup\suppset_{\MSP}}$, we have
    \begin{equation}
        \P\braces*{\monomial_S = 1 ~\vline~ \bX_{i,-\traversal^{(a)}\cup\suppset_{\MSP}} = \bx_{i,-\traversal^{(a)}\cup\suppset_{\MSP}}} = \frac{1}{2}.
    \end{equation}
    Combining this with the fact that the covariates are drawn uniformly from $\braces{\pm 1}^{\nfeats-1}$ as well as the symmetry of the splitting criterion, we see that the distribution of $\data^{(a)}$ is invariant to permutation of the covariate indices in $\coordindices\backslash\paren*{\traversal^{(a)}\cup\suppset_{\MSP}}$.
    As such, $J(\bx;\data^{(a)},\Theta)\backslash\suppset_{\MSP}$ is the prefix of a uniformly random permutation of $\coordindices\backslash\paren*{\traversal^{(a)}\cup\suppset_{\MSP}}$.

    \textit{Step 3: Bounding covariate selection in coupled paths.}
    For any $b \in \coordindices\backslash\paren*{\traversal^{(a)}\cup\suppset_{\MSP}}$, we use the previous step to compute
    \begin{equation}
        P\braces*{b \in J(\bx;\data^{(a)},\Theta)~\vline~\abs{J(\bx;\data^{(a)},\Theta)}} \leq \frac{\abs{J(\bx;\data^{(a)},\Theta)}}{\nfeats - \sparsity_{\MSP}-\traversalsize}.
    \end{equation}
    Taking a further expectation and using Lemma \ref{lem:bound_for_tree_depth} gives
    \begin{equation}
    \begin{split}
        \P\braces*{b \in J(\bx;\data^{(a)},\Theta)} & = \E\braces*{\P\braces*{b \in J(\bx;\data^{(a)},\Theta)~\vline~\abs{J(\bx;\data^{(a)},\Theta)}}} \\
        & \leq \frac{\log_2 \nsamples + 2}{\nfeats - \sparsity_{\MSP}-\traversalsize}.
    \end{split}
    \end{equation}

    \textit{Step 4: Decoupling implies covariate selection.}
    We claim the following inclusion of events:
    \begin{equation}
    \label{eq:non_msp_decoupling_covariate_selection}
        \bigcup_{a=1}^{\traversalsize}\braces*{J(\bx;\data,\Theta) \neq J(\bx;\data^{(a)},\Theta)} \subset \bigcup_{a=1}^{\traversalsize}\braces*{\sparsity_{\MSP} + a \in J(\bx;\data^{(a)},\Theta)}.
    \end{equation}
    To prove this, assuming the left hand side, let $t$ be the smallest index for which $k_t(\bx;\data^{(a)},\Theta) \neq k_t(\bx;\data,\Theta)$ for some $a$, which we assume without loss of generality is $a=1$.
    This implies in particular that the depth $t$ node along all $\traversalsize$ query paths are equal.
    Label this node as $\cell_t$.
    We now examine the possible values for $k_t(\bx;\data,\Theta)$.
    For $j \in \coordindices\backslash\paren*{\cellfeatindices(\cell_t)\cup\traversal^{(1)}}$, we have
    \begin{equation}
        \mathcal{O}(j,\cell_t,\data) = \mathcal{F}\paren*{\braces*{(X_{ij},Y_i)}_{i \in I(\cell)}} = \mathcal{O}(j,\cell_t,\data^{(1)}).
    \end{equation}
    Recall that $k_t(\bx;\data,\Theta)$ is defined to be the maximum of the left hand side over $j \notin \cellfeatindices(\cell_t)$.
    Meanwhile, $k_t(\bx;\data^{(1)},\Theta)$ is defined to be the maximum of the right hand side over $j \notin \cellfeatindices(\cell_t)\cup\traversal^{(1)}$.
    As such, we see that the only way in which these can be different values is if $b\coloneqq k_t(\bx;\data,\Theta) \in \traversal^{(1)}$.
    This would imply, however, that $k_t(\bx;\data^{(b-\sparsity_{\MSP})},\Theta) = b$ as we wanted.

    \textit{Step 5: Completing the proof.}
    Suppose $J(\bx;\data,\Theta)\cap \traversal \neq \emptyset$.
    If coupling does not hold (i.e., the left hand side of \eqref{eq:non_msp_decoupling_covariate_selection} is true), then by that same equation, we see that $\sparsity_{\MSP} + a \in J(\bx;\data^{(a)},\Theta)$ for some $a=1,2,\ldots,\traversalsize$.
    On the other hand if coupling does hold, then picking $\sparsity_{\MSP} + a \in J(\bx;\data,\Theta)\cap \traversal$, we have $\sparsity_{\MSP} + a \in J(\bx;\data,\Theta) = J(\bx;\data^{(a)},\Theta)$.
    In conclusion, we have
    \begin{equation}
    \begin{split}
        \P_{\data}\braces*{J(\bx;\data,\Theta)\cap \traversal \neq \emptyset } & \leq \sum_{a=1}^{\traversalsize} \P\braces*{\sparsity_{\MSP} + a \in J(\bx;\data^{(a)},\Theta)} \\
        & \leq \frac{\traversalsize(\log_2 \nsamples + 2)}{\nfeats - \sparsity_{\MSP}-\traversalsize+1}. \qedhere
    \end{split}
    \end{equation}
\end{proof}

\end{document}
